% 来源及取材范围见分题文件；此为整理者转录的正文，不是下载原件。
\begin{theorem}[Regular points of a zero set]\label{regptsdense}
Suppose $U\subset\C^n$ is a domain, $f\in\sO(U)$, and $f$ is not identically zero. Let $N=f^{-1}(0)$. Then there exists an open and dense (subspace topology) subset $N_{\mathit{reg}}\subset N$ such that at each $p\in N_{\mathit{reg}}$, after possibly reordering variables, $N$ can be locally (that is, in some neighborhood) written as
\[z_n=g(z_1,\ldots,z_{n-1})\]
for a holomorphic function $g$.
\end{theorem}
\begin{proof}
If $N$ is locally a graph at $p$, then it is a graph for every point of $N$ near $p$. So $N_{\mathit{reg}}$ is open. If for every point $p_0\in N$ and every neighborhood $W$ of $p_0$, we show that $N\cap W$ has a regular point, then $N_{\mathit{reg}}$ is dense. Replacing $N$ with $N\cap W$, it thus suffices to show $N_{\mathit{reg}}$ is nonempty.

Since $f$ is not identically zero, not all derivatives (of arbitrary order) of $f$ vanish identically on $N$. If some first order derivative of $f$ does not vanish identically on $N$, let $h=f$. Otherwise, suppose $k$ is such that a derivative of $f$ of order $k$ does not vanish identically on $N$, and all derivatives of $f$ of order less than $k$ vanish identically on $N$. Let $h$ be one of the derivatives of order $k-1$. We obtain a function $h\colon U\to\C$, holomorphic, vanishing on $N$, and such that without loss of generality the $z_n$ derivative does not vanish identically on $N$. Then there is some point $p\in N$ such that $\frac{\partial h}{\partial z_n}(p)\ne0$. We apply the implicit function theorem at $p$ to find $g$ such that
\[h\bigl(z_1,\ldots,z_{n-1},g(z_1,\ldots,z_{n-1})\bigr)=0,\]
and $z_n=g(z_1,\ldots,z_{n-1})$ is the unique solution to $h=0$ near $p$.

The zero set of $h$ contains $N$, the zero set of $f$. We must show equality near $p$. That is, we need to show that near $p$, every zero of $h$ is also a zero of $f$. Write $p=(p',p_n)$. The function $\xi\mapsto f(p',\xi)$ has an isolated zero in a small disc $\Delta$ around $p_n$ and is nonzero on the circle $\partial\Delta$. By Rouch\'e's theorem, $\xi\mapsto f(z',\xi)$ has a zero in $\Delta$ for all $z'$ sufficiently close to $p'$ (close enough to make $\sabs{f(p',\xi)-f(z',\xi)}<\sabs{f(p',\xi)}$ for all $\xi\in\partial\Delta$). Since $g(z')$ is the unique solution $z_n$ to $h(z',z_n)=0$ near $p$ and the zero set of $f$ is contained in the zero set of $h$, we are done.
\end{proof}
\begin{theorem}[Injective holomorphic mappings]
Suppose $U\subset\C^n$ is an open set and $f\colon U\to\C^n$ is holomorphic and one-to-one. Then the Jacobian determinant is never equal to zero on $U$.
In particular, if a holomorphic map $f\colon U\to V$ is one-to-one and onto for two open sets $U,V\subset\C^n$, then $f$ is biholomorphic.
\end{theorem}
The function $f$ is locally biholomorphic, in particular $f^{-1}$ is holomorphic, on the set where the Jacobian determinant
\[J_f(z)=\det Df(z)=\det\left[\frac{\partial f_k}{\partial z_\ell}(z)\right]_{k\ell}\]
is not zero. This follows from the inverse function theorem, which is just a special case of the implicit function theorem. The trick to prove the theorem above is to prove that $J_f$ is nowhere zero.

In one complex dimension, every nonconstant holomorphic function $f$ can, in the proper local holomorphic coordinates (and up to adding a constant), be written as $z^d$ for $d=1,2,\ldots$: Near a $z_0\in\C$, there exists a constant $c$ and a local biholomorphic $g$ with $g(z_0)=0$ such that $f(z)=c+\bigl(g(z)\bigr)^d$. So $f$ is one-to-one precisely if $d=1$.
\begin{proof}[Proof of the theorem]
We proceed by induction. We know the theorem for $n=1$. Suppose $n>1$ and suppose we know the theorem is true for dimension $n-1$.

Suppose for contradiction that $J_f=0$ somewhere. First suppose that $J_f$ is not identically zero. Find a regular point $q$ of the zero set of $J_f$. Write the zero set of $J_f$ near $q$ as
\[z_n=g(z_1,\ldots,z_{n-1})\]
for some holomorphic $g$. If we prove the theorem near $q$, we are done. Without loss of generality assume $q=0$. The biholomorphic (near the origin) map
\[\Psi(z_1,\ldots,z_n)=\bigl(z_1,z_2,\ldots,z_{n-1},z_n-g(z_1,\ldots,z_{n-1})\bigr)\]
takes the zero set of $J_f$ to the set given by $z_n=0$. Considering $f\circ\Psi^{-1}$ instead of $f$, we assume that $J_f=0$ on the set given by $z_n=0$. We also assume that $f(0)=0$.

If $J_f$ vanishes identically, then there is no need to do anything other than a translation. In either case, we may assume that $0\in U$, $f(0)=0$, and $J_f=0$ when $z_n=0$. Really, all we need is for the set where $J_f=0$ to be a sufficiently large set.

We wish to show that all the derivatives of $f$ in the $z_1,\ldots,z_{n-1}$ variables vanish whenever $z_n=0$. This would clearly contradict $f$ being one-to-one, as $f(z_1,\ldots,z_{n-1},0)$ would be constant. So for any point on $z_n=0$, consider one of the components of $f$ and one of the derivatives of that component. Without loss of generality, suppose the point is $0$, and for contradiction suppose $\frac{\partial f_1}{\partial z_1}(0)\ne0$. The map
\[G(z_1,\ldots,z_n)=\bigl(f_1(z),z_2,\ldots,z_n\bigr)\]
is biholomorphic on a small neighborhood of the origin. The function $f\circ G^{-1}$ is holomorphic and one-to-one on a small neighborhood. By the definition of $G$,
\[f\circ G^{-1}(w_1,\ldots,w_n)=\bigl(w_1,h(w)\bigr),\]
where $h$ is a holomorphic mapping taking a neighborhood of the origin in $\C^n$ to $\C^{n-1}$. The mapping
\[\varphi(w_2,\ldots,w_n)=h(0,w_2,\ldots,w_n)\]
is a one-to-one holomorphic mapping of a neighborhood of the origin in $\C^{n-1}$ to $\C^{n-1}$. By the induction hypothesis, the Jacobian determinant of $\varphi$ is nowhere zero.

If we differentiate $f\circ G^{-1}$, we notice $D(f\circ G^{-1})=Df\circ D(G^{-1})$. So at the origin
\[\det D(f\circ G^{-1})=(\det Df)(\det D(G^{-1}))=0.\]
We obtain a contradiction, as at the origin
\[\det D(f\circ G^{-1})=\det D\varphi\ne0.\qedhere\]
\end{proof}
